% ============================================================
% KAPITEL 7 — ZUSAMMENFASSUNG UND AUSBLICK
% ============================================================

\chapter{Zusammenfassung und Ausblick}
\label{ch:zusammenfassung}

% ------------------------------------------------------------
\section{Zusammenfassung}
\label{sec:zus-zusammenfassung}

Ausgangspunkt der Arbeit war die Frage, in welcher Form einem lokal betriebenen KI-Agenten eine 3D-Szene bereitgestellt werden sollte, damit er sie über natürliche Sprache manipulieren kann.
Zur Auswahl standen ein struktureller und ein visueller Weg, die beide eigene Stärken haben und jeweils einen Bereich abdecken, den der andere nicht erreicht.

Um diese Frage zu beantworten, wurde ein prototypisches System entwickelt, das dieselbe Szene strukturell, visuell und hybrid bereitstellt und über ein gemeinsames Interface manipulierbar macht.
Die Repräsentationsformen wurden unter gleichen Bedingungen an dreizehn kontrolliert aufgebauten Szenarien verglichen, wobei jede Kombination aus Szenario und Variante in 30 Läufen gemessen und der Erfolg aus dem tatsächlichen Zustand der Szene bestimmt wurde.

Welche Repräsentation am besten abschneidet, hing vom Referenztyp ab, und eine naive Kombination beider Schichten vereinte deren Stärken nicht automatisch.
Erst eine Instruktion, die den Schichten getrennte Zuständigkeiten zuweist, brachte beide auf dem untersuchten Szenarioset zusammen.

Für den Demonstrator wurde diese kompetenzgetrennte Variante in die immersive Zielumgebung übertragen.
Er zeigt, dass die Kette vom Befehl bis zur sichtbaren Manipulation in der Brille funktioniert.
Per Video dokumentiert ist die Fassung mit Bedienmenü, die später ergänzte Spracheingabe nicht.


% ------------------------------------------------------------
\section{Beantwortung der Forschungsfragen}
\label{sec:zus-rq}

\subsection*{Erste Forschungsfrage}

\begin{quote}
\textit{Welchen Einfluss hat die Form der Szenenrepräsentation, strukturell, visuell oder hybrid, auf den Ausführungserfolg der natürlichsprachlichen Szenenmanipulation durch einen lokal betriebenen Agenten, und zwar aufgeschlüsselt nach der Art der sprachlichen Referenz?}
\end{quote}

Für die erste Forschungsfrage lautet das Fazit, dass die Repräsentationsform einen großen Einfluss auf den Ausführungserfolg hat und dieser Einfluss vom Referenztyp abhängt.
Referenzen, deren Merkmal in den Daten kodiert ist, sowie Referenzen auf verdeckte Objekte werden von der strukturellen Form getragen, die als einzige ein im Bild verdecktes Objekt ansprechen kann.
Referenzen auf rein visuelle Eigenschaften wie ein Oberflächenmuster lassen sich nur über die visuelle Form auflösen.

Die naive hybride Form übernimmt den Vorteil der strukturellen Form, den der visuellen nicht zuverlässig.
In der perzeptuellen Kategorie erreicht sie nur 10,7 von 20 möglichen Erfolgen, die visuelle Form alle 20.
Eine Kombination ist also nicht automatisch so gut wie die bessere der beiden Einzelformen.

\subsection*{Zweite Forschungsfrage}

\begin{quote}
\textit{Lässt sich das Zusammenspiel der strukturellen und der visuellen Schicht allein durch die Gestaltung der Instruktion so steuern, dass die hybride Repräsentation die Stärken beider Schichten vereint?}
\end{quote}

Für das untersuchte Szenarioset lässt sich diese Frage mit Ja beantworten, ohne dass die Modelle verändert werden müssen.
Eine Priorisierung einer Schicht reicht dafür nicht, sie verschiebt nur das Gewicht und tauscht eine Kompetenz gegen eine andere.
Eine klare Trennung der Zuständigkeiten, bei der die visuelle Beschreibung das gemeinte Objekt bestimmt und die strukturelle Schicht alle geometrischen Werte liefert, bringt beide Schichten dagegen zusammen.
Die kompetenzgetrennte Variante erreicht so in der perzeptuellen Kategorie 19,7 von 20 und behält gleichzeitig die Stärke der strukturellen Form.


% ------------------------------------------------------------
\section[Beitrag]{Wissenschaftlicher Beitrag}
\label{sec:zus-beitrag}

Inhaltlich zeigt die Arbeit an einem kontrollierten Fall, dass sich strukturelle und visuelle Information nicht automatisch ergänzen.
Ob eine Kombination beide Informationsarten sinnvoll zusammenführt oder ob eine die andere verdrängt, hängt davon ab, wie die Kombination gestaltet ist.
Gewonnen wurde dieser Befund mit einem Modellpaar an dreizehn Szenarien.
Er legt nahe, die Kombination mehrerer Informationsquellen auch bei anderen Agenten als Gestaltungsaufgabe zu sehen und nicht einfach als zusätzliche Information.

Methodisch führt die Arbeit eine nach Referenztypen aufgeschlüsselte Auswertung ein, bei der der Erfolg für jeden Referenztyp einzeln gemessen wird und nicht als ein Gesamtwert.
Erst dadurch werden die exklusiven Zuständigkeiten der Repräsentationen sichtbar, die bei einem pauschalen Vergleich verborgen geblieben wären.
Da dieses Schema nicht an das konkrete System gebunden ist, sollte es sich auch auf ähnliche Vergleiche übertragen lassen.

Praktisch zeigt die Arbeit, dass sich die gesamte Kette aus 3D-Umgebung, Vermittlungsschicht, Agent und zwei gleichzeitig betriebenen lokalen Modellen umsetzen und bis in die immersive Zielumgebung betreiben lässt.


% ------------------------------------------------------------
\section{Ausblick}
\label{sec:zus-ausblick}

Die meisten Ansatzpunkte für weitere Arbeiten ergeben sich direkt aus den Einschränkungen in Abschnitt~\ref{sec:eval-validity}.
Als Erstes sollte geprüft werden, ob die kompetenzgetrennte Variante ihren Vorteil auch außerhalb des Szenariosets behält, auf dem sie entwickelt wurde.
Dafür bräuchte es neue Szenarien, vor allem mit anderen perzeptuellen Eigenschaften als einem Streifenmuster und mit anderen Formulierungen für dieselbe Referenz, sowie eine Fassung der Instruktion ohne \enquote{striped} und \enquote{most salient}, in der die Trennung der Zuständigkeiten rein allgemein formuliert ist.
Erst wenn eine solche Instruktion auch auf neuen Szenarien denselben Vorteil zeigt, ließe sich von einem übertragbaren Gestaltungsprinzip sprechen.

Ähnlich wichtig ist eine Messung, die den Reihenfolgeeffekt in Kategorie D ausschließt.
Dafür würde es reichen, die entscheidenden Szenarien mit vertauschter Objektreihenfolge zu wiederholen und die Reihenfolge der beiden Schichten im Kontext unabhängig von der Instruktion zu verändern.
Beides ist ohne neuen Aufbau möglich.
Bei jeder Wiederholung sollten zudem Werkzeugaufrufe und Endpositionen protokolliert werden, womit sich die offenen Fragen zu C2 und C4 beantworten ließen.

Weitere Ansatzpunkte betreffen die Effizienz und den Umfang des Systems.
Die visuelle Schicht abzurufen dauert spürbar länger als der Abruf des strukturellen Zustands, da ein zusätzlicher Inferenzschritt des Sprach-Bild-Modells nötig ist.
Gemessen wurde dieser Unterschied nicht.
Eine systematische Erhebung von Latenz, Schrittzahl und Rechenaufwand würde diesen bisher nur qualitativ bekannten Trade-off greifbar machen, was vor allem für eine interaktive immersive Anwendung wichtig ist.
Mit steigender Objektzahl dürfte vor allem die visuelle Schicht an ihre Grenzen kommen, weil die Markierung vieler Objekte schnell unübersichtlich wird.
Die Menge der Operationen ließe sich um Drehungen und Skalierungen erweitern sowie um Befehle mit mehreren Manipulationen, deren eigene Schwierigkeiten sich hier bereits angedeutet haben.
Offen sind auch blickpunktabhängige Referenzen, die sich auf die Sicht des Betrachters beziehen und gerade für immersive Umgebungen typisch sind.
Jüngere Arbeiten zur blickpunktabhängigen Segmentierung zeigen, dass solche Bezüge ohne einen expliziten Blickpunkt mehrdeutig sind und von einer ausdrücklichen Darstellung des Blickwinkels profitieren~\cite{nanri2026}.
Eine saubere Behandlung dieses Referenztyps wäre deshalb eine sinnvolle Erweiterung der Taxonomie.

Auf wissenschaftlicher Ebene ließe sich prüfen, ob die gefundenen Kompetenzprofile und die Grenzen der naiven Kombination auch für größere oder anders aufgebaute Modelle gelten oder ob ein größeres Modell das beobachtete Problem von selbst löst.
Ein gezieltes Experiment könnte außerdem die drei offenen Erklärungen voneinander trennen, indem Formulierung, visuelle Eindeutigkeit und Reihenfolge der Objekte unabhängig voneinander verändert werden.
Interessant wäre auch ein Vergleich der zweistufigen Pipeline mit einem einzelnen multimodalen Modell, dem Bild und strukturelle Daten direkt vorliegen.
Damit ließe sich zeigen, ob der beobachtete Effekt von der Pipeline abhängt oder auch unabhängig davon auftritt.
Ob sich das Prinzip der Kompetenztrennung auch auf andere multimodale Aufgaben übertragen lässt, kann Gegenstand zukünftiger Arbeiten werden.

Aus Sicht der Anwendung fehlt noch eine Erprobung mit echten Nutzern.
Die Bewertung in dieser Arbeit ist rein systembasiert, und eine qualitative Nutzerstudie im Demonstrator könnte zeigen, wie Menschen die sprachgesteuerte Manipulation in der immersiven Umgebung tatsächlich erleben.
